% from ega4-18-fr.tex
\oldpage[IV-4]{218}

En particulier, on peut donc écrire
\[
  g=t+g_2
\]
avec $t\in\mathfrak J^h(R\cap\mathfrak R^dN)$ et $g_2\in\mathfrak R^{d+1}N+\mathfrak R^d\mathfrak L^hN$; en outre, comme $t\in R$ et $g\in R$, on a $g_2\in R$. D'autre part, utilisant l'hypothèse $h\geq m-d+1$, on voit que $t\in\mathfrak L^{m+1}N$. L'élément $g_2$ répond donc à toutes les conditions énoncées au début de (19.7.1.8), ce qui achève la démonstration de (19.7.1).

\begin{remarks}[19.7.2]
\label{IV.19.7.2}
\begin{enumerate}
  \item[(i)] Lorsque l'on suppose que $\operatorname{gr}_{\mathfrak R}^*(A)$ est un $(A/\mathfrak R)$-module \emph{plat}, l'hypothèse que la suite $\mathbf f$ est $(A/\mathfrak R)$-régulière entraîne qu'elle est aussi $\operatorname{gr}_{\mathfrak R}^*(A)$-régulière ($0$, 15.1.14). On notera que ce sera le cas lorsque $A$ et $A/\mathfrak R$ sont des anneaux réguliers ($0$, 17.3.6): en effet, $\mathfrak R$ est alors un idéal régulier de $A$ (19.1.2), donc quasi-régulier, et en localisant aux idéaux maximaux de $A$ contenant $\mathfrak R$ et utilisant ($0$, 15.1.7), on voit que $\operatorname{gr}_{\mathfrak R}^*(A)$ est un $(A/\mathfrak R)$-module \emph{projectif}.
  \item[(ii)] Il serait intéressant de pouvoir prouver la dernière conclusion de (19.7.1) en supposant seulement qu'il existe une $A$-algèbre $B$ qui est un anneau noethérien, pour lequel $\mathfrak JB$ est contenu dans le radical de $B$, et que $M$ est un $B$-module de type fini.
\end{enumerate}
\end{remarks}

\begin{corollary}[19.7.3]
\label{IV.19.7.3}
Soient $A$ un anneau local noethérien, $\mathfrak m$ son idéal maximal, $\mathfrak R$ un idéal de $A$ tel que l'anneau $A/\mathfrak R$ soit régulier et de dimension $n$, $M$ un $A$-module de type fini. Les conditions suivantes sont équivalentes:
\begin{enumerate}
  \item[a)] $\operatorname{gr}_{\mathfrak R}^*(M)$ est un $(A/\mathfrak R)$-module plat.
  \item[b)] Si $P(T)$ et $Q(T)$ sont les séries de Poincaré des $(A/\mathfrak m)$-modules gradués $\operatorname{gr}_{\mathfrak m}^*(M)$ et $\operatorname{gr}_{\mathfrak R}^*(M)\otimes_{A/\mathfrak R}(A/\mathfrak m)$, on a
\end{enumerate}
\begin{equation}
  Q(T)=P(T)(1-T)^n.
  \tag{19.7.3.1}\label{IV.19.7.3.1}
\end{equation}

Soit $\mathbf f=(f_1,\ldots,f_n)$ une suite d'éléments de $A$ tels que les images des $f_i$ dans $A/\mathfrak R$ forment un système régulier de paramètres de $A/\mathfrak R$ ($0$, 17.1.6), de sorte que $\mathbf f$ est une suite $(A/\mathfrak R)$-régulière; en outre, si $\mathfrak J$ est l'idéal engendré par $\mathbf f$, on a $\mathfrak J+\mathfrak R=\mathfrak m$, et comme $A/\mathfrak m$ est un corps, $\operatorname{gr}_{\mathfrak m}^*(M)$ est un $(A/\mathfrak m)$-module plat. Comme $A$ est noethérien et $M$ un $A$-module de type fini, on peut appliquer l'équivalence de $a)$ et $b)$ dans (19.7.1). En outre, comme $A/\mathfrak m$ est un corps, et que $\psi_M$ est surjectif, le fait que $\psi_M$ soit bijectif équivaut à dire que pour tout $h\geq0$, $\operatorname{gr}_{\mathfrak m}^h(M)$ et le sous-module des éléments de degré $h$ de $(\operatorname{gr}_{\mathfrak R}^*(M)\otimes_{A/\mathfrak R}(A/\mathfrak m))[T_1,\ldots,T_n]$ ont même rang sur $A/\mathfrak m$; mais pour le second de ces modules le rang en question est évidemment égal à
\[
  \sum_{r=0}^h\binom{r+n-1}{n-1}\operatorname{rg}(\operatorname{gr}_{\mathfrak R}^{h-r}(M)\otimes_{A/\mathfrak R}(A/\mathfrak m)).
\]
Comme
\[
  (1-T)^{-n}=\sum_{r=0}^{\infty}\binom{r+n-1}{n-1}T^r,
\]
cela prouve que la relation (19.7.3.1) est nécessaire et suffisante pour que $\psi_M$ soit bijectif.
\end{corollary}

\begin{corollary}[19.7.4]
\label{IV.19.7.4}
Soient $X$ un préschéma localement noethérien, $Y$ un sous-préschéma fermé de $X$, régulier et connexe, $\sh F$ un $\sh O_X$-Module cohérent normalement plat (6.10.1) le long de $Y$. Alors il existe une série formelle $R\in\mathbb{Q}[[T]]$ (indépendante de $x\in Y$) telle que, pour tout $x\in Y$, la série de Poincaré de $\operatorname{gr}_{\mathfrak m_x}^*(\sh F_x)$ soit donnée par la formule
\begin{equation}
  P_x(\sh F)(T)=R(T)(1-T)^{-n}\qquad\text{où }n=\dim(\sh O_{Y,x}).
  \tag{19.7.4.1}\label{IV.19.7.4.1}
\end{equation}

